\section{Implications for stateful agents}
\label{sec:C03:implications}
The foundation model is a conditional computation embedded in a stateful control system. Its capabilities remain important: stronger representations, broader training coverage, and better action distributions can reduce the burden on the harness. Nevertheless, the useful unit of analysis is the composition of tokenization, conditioning, decoding, interpretation, and execution. A model can only act on information available through that composition. Persistent storage helps only when the evidence reaches a representation and interface through which it can affect the correct decision.

This observation constrains claims about memory and self-improvement. Replacing a tokenizer or changing special-message serialization can alter the meaning of a stored token sequence even if the original text survives. Updating model weights can change how a procedural memory is interpreted. A context-editing harness can invalidate cached computational state. Compatibility is therefore more than a database schema version: it includes the model, tokenizer, prompt construction, action grammar, tool semantics, and authorization epoch. The detailed consistency analysis in Chapter~\ref{ch:C17} builds on these interfaces rather than treating memory as an isolated retrieval component.

The engineering implication is to maintain explicit boundaries while studying their interactions. Save canonical external records in a representation that preserves their intended semantics and provenance, reconstruct model-specific inputs under a declared policy, and evaluate behavior after every material interface change. Do not infer indefinite competence from a larger window, universal memory from a recurrent state, or safe execution from valid JSON. Each property has a measurable operational definition, a specific implementation boundary, and failure conditions that can be tested independently before assessing the complete agent.

\paragraph{A precise limitation of compressed architectural state.}
The state-aliasing argument can be strengthened into a decision bound. Suppose two histories $h_0,h_1$ occur with probabilities $p$ and $1-p$, produce the same exposed state $z$, and require distinct unique correct actions $a_0,a_1$. Assume no additional observation arrives before the action. Any randomized policy conditioned only on $z$ has success probability
\begin{equation}
 p\pi(a_0\mid z)+(1-p)\pi(a_1\mid z)
 \leq \max\{p,1-p\}.
\end{equation}
The proof uses $\pi(a_0\mid z)+\pi(a_1\mid z)\leq1$: the linear objective is maximized by allocating all probability to the more frequent required action. Hence the minimum error is at least $\min\{p,1-p\}$. This is an elementary deduction under the stated observational restriction, not a new impossibility result for any particular architecture. External retrieval can remove the restriction by supplying evidence that distinguishes the histories. A permission check can prevent execution of a wrong action, but cannot reconstruct the missing task information by itself.

This bound suggests a targeted architectural evaluation. Construct paired histories that differ in one decision-critical fact, then measure whether the system's retained representation and available retrieval preserve that distinction. Keep the final query identical. Include an explicit-record baseline, a recurrent-state-only baseline where accessible, and an external-memory baseline with equal retrieval budgets. Balance the two histories so that a default action cannot obtain a misleadingly high score through prior frequency. The test measures a particular lost distinction; it cannot establish that every important distinction survives. Nevertheless, it is more diagnostic than averaging many easy tasks whose correct action is unchanged by the omitted history.
